% Source: https://artofproblemsolving.com/community/c3h1944534_mock_amc_8a_final_five
% Stable identity: aops-topic-1944534-q21
First work modulo $5$: $n^2\equiv-1\equiv4\pmod5$ forces $n\equiv2$ or $3\pmod5$.
\begin{itemize}
\item $n=5k+2$: $n^2+1=25k^2+20k+5$, which is divisible by $25$ iff $20k+5\equiv0\pmod{25}$, i.e. $4k+1\equiv0\pmod5$, i.e. $k\equiv1\pmod5$, so $n\equiv7\pmod{25}$.
\item $n=5k+3$: $n^2+1=25k^2+30k+10$, divisible by $25$ iff $6k+2\equiv0\pmod5$, i.e. $k\equiv3\pmod5$, so $n\equiv18\pmod{25}$.
\end{itemize}
So exactly the residues $7$ and $18$ modulo $25$ work. The range $0\le n\le1999$ consists of $80$ complete blocks of $25$, giving $160$ values, and $2000\le n\le2019$ covers residues $0$ to $19$, which include $7$ and $18$: $2$ more. The answer is $\boxed{\textbf{(D) }162}$.
